\section{A Sketch of the Proof}\label{sec:overview}

We first use flag algebras to prove the bound $\frac{5}{4}\Delta(G)^2$
for $\omega(L(G)^2)$ in simple graphs. By counting parallel edges
with their multiplicities, we extend the bound to loopless multigraphs.
This settles the even-degree case.
For odd $\Delta=\Delta(G)\geq 9$, we show that if
$\omega(L(G)^2)>\Delta^2+\lfloor\frac{\Delta}{2}\rfloor^2$,
then $L(G)^2$ has a clique exceeding this bound whose vertices
correspond to a set of edges of $G$ that includes both a pair of
parallel edges and five edges forming an induced five-cycle in the
underlying simple graph of $G$. Counting around such cycles gives a
bound that contradicts the weighted flag inequalities. We handle the remaining small odd degrees separately,
thereby completing the proof.

\subsection{The Classical Flag Algebra Method}\label{subsec:flag-algebra}

Let $G$ be a graph. For a graph $H$ with $|V(H)|\leq |V(G)|$,
choose a subset $U\subseteq V(G)$ of size $|V(H)|$ uniformly at
random. The \emph{induced subgraph density} of $H$ in $G$ is the
probability that $G[U]$ is isomorphic to $H$. A \emph{type}
$\sigma$ is a graph with a distinct label on each vertex.
A \emph{$\sigma$-flag} is a partially labeled graph whose labeled
vertices induce $\sigma$ with the same labels. When the type
$\sigma$ is clear from the context, we call a $\sigma$-flag
simply a flag. Two flags are considered the same if there is an
isomorphism between them that preserves the labels. Fix a
labeled induced copy of $\sigma$ in $G$. For a $\sigma$-flag $F$
with $|V(F)|\leq |V(G)|$, choose uniformly at random a set of
$|V(F)|-|V(\sigma)|$ vertices outside this copy. The \emph{flag
density} of $F$ in $G$ is the probability that the chosen vertices
together with the fixed vertices induce a copy of $F$ with the
labels preserved. Fix a type $\sigma$. Let $\mathcal F^\sigma$
be the set of $\sigma$-flags, and let $\mathcal F^\sigma_n$ be
the set of $\sigma$-flags on $n$ vertices. Consider formal sums
$\sum_{J\in\mathcal F^\sigma}a_JJ$ with real coefficients, only
finitely many of which are nonzero. We add two such sums by
adding their coefficients at each flag $J$. Multiplication by
a scalar $\lambda$ sends each coefficient $a_J$ to $\lambda a_J$.
With these operations, it is easy to see that the formal sums
form a real vector space with basis $\mathcal F^\sigma$. We
denote this space by $\mathbb R\mathcal F^\sigma$.

For flags $F$ and $J$ with $|V(F)|\leq |V(J)|$, write $p(F;J)$
for the flag density of $F$ in $J$. We use $p(F;G)$ in the same
sense when the labeled copy of $\sigma$ in $G$ has been fixed.
For each flag $F$ and each
integer $n\geq |V(F)|$, form the vector
\[
  F-\sum_{J\in\mathcal F^\sigma_n}p(F;J)J.
\]
Let $\mathcal K^\sigma$ be the subspace of
$\mathbb R\mathcal F^\sigma$ spanned by these vectors as $F$
and $n$ vary. Put
$\mathcal A^\sigma=\mathbb R\mathcal F^\sigma/\mathcal K^\sigma$.
Two formal sums represent the same element of $\mathcal A^\sigma$
if their difference lies in $\mathcal K^\sigma$.

To define multiplication, let $F_1$ and $F_2$ be flags with
$r_1$ and $r_2$ unlabeled vertices, respectively. Put
$n=|V(\sigma)|+r_1+r_2$. For each $J\in\mathcal F^\sigma_n$,
choose $r_1$ of its unlabeled vertices uniformly at random
and put the remaining $r_2$ in the second group. Each group,
together with the labeled vertices of $J$, induces a flag.
Let $p(F_1,F_2;J)$ be the probability that the chosen group
gives $F_1$ and the remaining group gives $F_2$. The product
$F_1\cdot F_2$ is the element of $\mathcal A^\sigma$ represented by
\[
  \sum_{J\in\mathcal F^\sigma_n}p(F_1,F_2;J)J.
\]
To multiply two formal sums $\sum_i a_iF_i$ and $\sum_j b_jH_j$,
take the flag product $F_i\cdot H_j$ for each pair of terms, one
from each sum. Multiply each such product by the corresponding
coefficient $a_ib_j$ and add the results. Thus
\[
  \left(\sum_i a_iF_i\right)\cdot\left(\sum_j b_jH_j\right)
  =\sum_{i,j}a_ib_j(F_i\cdot H_j).
\]
Razborov showed that the product is well defined on the quotient
$\mathcal A^\sigma$~\cite{Razborov2007}. With this product,
$\mathcal A^\sigma$ is the \emph{flag algebra of type $\sigma$}.

Razborov's flag algebra method uses products of flag densities
to derive inequalities among induced subgraph
densities~\cite{Razborov2007,CarliSilvaOliveiraSato2016}.
Let $F_1$ and $F_2$ be $\sigma$-flags with $r_1$ and $r_2$
unlabeled vertices, respectively. Fix a labeled induced copy
of $\sigma$ in $G$, and suppose that the number of vertices
outside this copy is at least $\max\{1,r_1+r_2\}$. Choose
$r_1+r_2$ of these vertices uniformly at random. Together with
the labeled vertices, they induce a flag $J$. Choose $r_1$ of
the sampled vertices uniformly at random for the first group
and put the remaining $r_2$ in the second group. Given $J$,
the probability that the two groups, each taken together with
the labeled vertices, induce $F_1$ and $F_2$, respectively, is
$p(F_1,F_2;J)$. Since the probability of obtaining $J$ at the
first step is $p(J;G)$, the probability of obtaining both
$F_1$ and $F_2$ is
\[
  \sum_{J\in\mathcal F^\sigma_{|V(\sigma)|+r_1+r_2}}
  p(J;G)\,p(F_1,F_2;J).
\]
Now let $W$ be the set of vertices of $G$ outside the fixed labeled
copy of $\sigma$. Choose one of the $r_1$-element subsets of $W$
uniformly at random and call it $U_1$. Independently of this
choice, choose one of the $r_2$-element subsets of the same set
$W$ uniformly at random and call it $U_2$. Let $A$ be the event
that, for each $i\in\{1,2\}$, the subgraph induced by $U_i$ and
the labeled vertices gives the flag $F_i$. By independence,
$\Pr(A)=p(F_1;G)p(F_2;G)$.

Let $B$ be the event that $U_1\cap U_2=\varnothing$. By
independence and symmetry, conditional on $B$, all ordered
pairs $(S_1,S_2)$ of disjoint subsets of $W$ with $|S_1|=r_1$
and $|S_2|=r_2$ are equally likely. To count these pairs, choose
the union $S_1\cup S_2$ and then choose $S_1$ within it. There
are $\binom{|W|}{r_1+r_2}$ choices for the union and
$\binom{r_1+r_2}{r_1}$ choices for $S_1$ within each union.
Thus each pair has conditional probability
$\frac{1}{\binom{|W|}{r_1+r_2}\binom{r_1+r_2}{r_1}}$.

When the union is sampled first, each ordered pair $(S_1,S_2)$
of disjoint subsets of $W$ with $|S_1|=r_1$ and $|S_2|=r_2$
has probability
$\frac{1}{\binom{|W|}{r_1+r_2}\binom{r_1+r_2}{r_1}}$.
Under independent sampling, this is also the conditional
probability of $(U_1,U_2)=(S_1,S_2)$ given $B$. The two sampling
rules therefore give the same probability that the first set
and the second set, each together with the labeled vertices,
induce $F_1$ and $F_2$, respectively. For the experiment that
samples the union first, we computed this probability by
conditioning on the flag $J$ induced by the union and the
labeled vertices. Thus
\[
  \Pr(A\mid B)
  =\sum_{J\in\mathcal F^\sigma_{|V(\sigma)|+r_1+r_2}}
    p(J;G)\,p(F_1,F_2;J).
\]

It remains to compare $\Pr(A\mid B)$ with the probability $\Pr(A)$
under independent sampling. Since
$\Pr(A)=\Pr(A\mid B)\Pr(B)+\Pr(A\cap B^c)$, their difference is
bounded in absolute value by $\Pr(B^c)$. To bound $\Pr(B^c)$,
fix the outcome of $U_1$. The set $U_2$ is still a uniform
$r_2$-element subset of $W$, so each vertex of $U_1$ belongs
to $U_2$ with probability $\frac{r_2}{|W|}$. The event $B^c$
occurs if at least one of the $r_1$ vertices of $U_1$ belongs
to $U_2$. The union bound over the vertices of $U_1$ gives
$\Pr(B^c)\leq\frac{r_1r_2}{|W|}$, since the bound holds for every
choice of $U_1$. Since $\Pr(A)=p(F_1;G)p(F_2;G)$ and
$|W|=|V(G)|-|V(\sigma)|$, we have
\begin{equation}\label{eq:flag-product-estimate}
  \left|
    p(F_1;G)p(F_2;G)
    -\sum_{J\in\mathcal F^\sigma_{|V(\sigma)|+r_1+r_2}}
        p(F_1,F_2;J)p(J;G)
  \right|
  \leq \frac{r_1r_2}{|V(G)|-|V(\sigma)|}.
\end{equation}

Fix a list of $\sigma$-flags $F_1,\ldots,F_t$ and a real
symmetric positive semidefinite $t\times t$ matrix $Q=(Q_{ij})$.
For each $i$, let $r_i=|V(F_i)|-|V(\sigma)|$. Take $G$ large
enough for \eqref{eq:flag-product-estimate} to apply to every
pair $F_i,F_j$, and assume $|V(G)|\geq 2|V(\sigma)|$. Let
$\theta$ be an injective assignment of the labels of $\sigma$
to vertices of $G$ such that the assigned vertices, with
their labels, induce $\sigma$. Write $p_\theta(F;G)$ for
the density of a flag $F$ in $G$ with this assignment, and
let $v_\theta$ be the column vector with $i$th entry
$p_\theta(F_i;G)$. Since $Q$ is positive semidefinite,
\[
  0\leq v_\theta^{\mathsf T}Qv_\theta
    =\sum_{i,j=1}^{t}Q_{ij}\,
       p_\theta(F_i;G)p_\theta(F_j;G).
\]

For each pair $F_i,F_j$, expand the flag product $F_i\cdot F_j$
using flags on $|V(\sigma)|+r_i+r_j$ vertices, as in the definition
of multiplication. Multiply each expansion by $Q_{ij}$, add the
expansions, and collect terms to obtain the formal sum
$\sum_J c_JJ$. Thus the pair $F_i,F_j$ contributes
$Q_{ij}p(F_i,F_j;J)$ to $c_J$ when
$|V(J)|=|V(\sigma)|+r_i+r_j$. The coefficients $c_J$ depend only
on the fixed flags and $Q$.

By \eqref{eq:flag-product-estimate}, replacing
$p_\theta(F_i;G)p_\theta(F_j;G)$ by the expansion of
$F_i\cdot F_j$ with each flag replaced by its density in $G$
has error at most $\frac{r_ir_j}{|V(G)|-|V(\sigma)|}$.
Multiplying these bounds by $|Q_{ij}|$ and summing gives
\[
  \left|v_\theta^{\mathsf T}Qv_\theta
       -\sum_J c_Jp_\theta(J;G)\right|
  \leq
  \frac{\sum_{i,j=1}^{t}|Q_{ij}|r_ir_j}
       {|V(G)|-|V(\sigma)|}.
\]
Here and below, $J$ runs over the flags in the product
expansions. Since $|V(G)|\geq 2|V(\sigma)|$, the right-hand
side is at most $\frac{C}{|V(G)|}$, where
$C=2\sum_{i,j=1}^{t}|Q_{ij}|r_ir_j$. The nonnegativity of
$v_\theta^{\mathsf T}Qv_\theta$ now gives
\begin{equation}\label{eq:flag-density-lower-bound}
  \sum_J c_Jp_\theta(J;G)\geq-\frac{C}{|V(G)|}.
\end{equation}

Let $\Theta$ be the set of all injective assignments of the
labels of $\sigma$ to vertices of $G$. For each $\theta\in\Theta$
whose labeled vertices induce $\sigma$, use the densities
$p_\theta(J;G)$ defined above. For every other $\theta\in\Theta$,
put $p_\theta(J;G)=0$ for each $J$. The left-hand side of
\eqref{eq:flag-density-lower-bound} is then zero for those
assignments, so the inequality holds for every $\theta\in\Theta$.
Summing \eqref{eq:flag-density-lower-bound} over $\Theta$ and
dividing by $|\Theta|$ gives
\begin{equation}\label{eq:averaged-flag-density-lower-bound}
  \sum_J c_J
    \left(\frac{1}{|\Theta|}\sum_{\theta\in\Theta}p_\theta(J;G)\right)
  \geq -\frac{C}{|V(G)|}.
\end{equation}

Fix a flag $J$ in \eqref{eq:averaged-flag-density-lower-bound}, and
put $k=|V(J)|$, $s=|V(\sigma)|$ and $N=|V(G)|$. Choose
$\theta\in\Theta$ uniformly at random and then choose $k-s$ vertices
uniformly from the vertices that received no label. Let $S$ be
the union of the $k-s$ chosen vertices and the $s$ vertices
assigned labels by $\theta$. The two choices give $J$ if $G[S]$ with its assigned labels is
isomorphic to $J$. For a fixed $\theta$, the probability of
obtaining $J$ is $p_\theta(J;G)$. This also holds
under the zero convention when the labeled vertices do not
induce $\sigma$. When $\theta$ is chosen uniformly from $\Theta$,
the probability of obtaining $J$ in the two steps is
\[
  \frac{1}{|\Theta|}\sum_{\theta\in\Theta}p_\theta(J;G).
\]
We call this probability the \emph{average flag density} of $J$
in $G$.

For a fixed $k$-element set $S$ and a fixed injective assignment
$\theta$ of the labels within $S$, the probability of obtaining
the pair $(S,\theta)$ by assigning the labels first is
$\frac{1}{s!\binom{N}{s}\binom{N-s}{k-s}}$. Indeed, we must first
choose $\theta$ from the $s!\binom{N}{s}$ injective assignments
in $G$ and then choose the $k-s$ unlabeled vertices of $S$.
If we choose $S$ uniformly first and then choose an injective
assignment within $S$ uniformly, the probability of the same pair
is $\frac{1}{\binom{N}{k}s!\binom{k}{s}}$. These probabilities are
equal by the identity
$\binom{N}{s}\binom{N-s}{k-s}=\binom{N}{k}\binom{k}{s}$,
so the two sampling procedures give the same distribution.

Write $H_J$ for the unlabeled graph on the $k$ vertices of the
fixed flag $J$, with the same edges as $J$. For an unlabeled
graph $H$, write $p(H;G)$ for its induced subgraph density in $G$.
When $S$ is chosen first, the graph on which we assign labels is
the $k$-vertex induced subgraph $G[S]$ of the host $G$. A label
assignment within $S$ can give the flag $J$ only when
$G[S]\cong H_J$, an event of probability $p(H_J;G)$.
Choose an injective assignment of the labels of $\sigma$ to
vertices of $H_J$ uniformly at random. Let $q_\sigma(J)$ be the
probability that this assignment gives the flag $J$. Given a set
$S$ with $G[S]\cong H_J$, a uniform injective assignment of the
labels within $S$ gives $J$ with probability $q_\sigma(J)$. Hence
the average flag density of $J$ in $G$ satisfies
\[
  \frac{1}{|\Theta|}\sum_{\theta\in\Theta}p_\theta(J;G)
    =q_\sigma(J)p(H_J;G).
\]
Substituting this average into
\eqref{eq:averaged-flag-density-lower-bound} gives
\begin{equation}\label{eq:flag-unlabeled-lower-bound}
  \sum_J c_Jq_\sigma(J)p(H_J;G)
    \geq -\frac{C}{|V(G)|}.
\end{equation}

Inequality \eqref{eq:flag-unlabeled-lower-bound} gives a lower
bound on a linear combination of induced subgraph densities.
Flag algebra inequalities can be used to bound the induced
subgraph density of a fixed graph in a given graph class.
Let $d(G)$ be the density we want to bound and let $U$ be a
proposed upper bound. Suppose
identities for the graph class allow us to write $U-d(G)$ as
the left-hand side of \eqref{eq:flag-unlabeled-lower-bound}
plus a nonnegative linear combination of subgraph densities. Then
$U-d(G)\geq-\frac{C}{|V(G)|}$, or
$d(G)\leq U+\frac{C}{|V(G)|}$. The error tends to zero as
$|V(G)|$ grows. For example, when $d(G)$ is the edge density
in a triangle-free graph, the identity $p(K_3;G)=0$ can be
used in finding such an expression for $U-d(G)$.

\subsection{The Flag Algebra Method in Our Proof}\label{subsec:adapted-flags}

We adapt the local flag algebra method of Davey, Hurley, de Joannis de
Verclos, Kang and Volec~\cite{DaveyEtAlLocal2026,DaveyEtAlStrong2026},
who use the maximum degree to normalize counts of small graphs.
In our problem, we wish to express the size of a given strong clique
in terms of such counts. We mark the edges of the clique and record
both the marked and the unmarked edges, since an unmarked edge may
join two disjoint marked edges, as in \textbf{Figure~\ref{fig:marked-maps}}(a).

Suppose first that $G$ is simple, has maximum degree $\Delta$, and
contains a strong clique of $m>0$ edges. By
\textbf{Lemma~\ref{lem:regularization}}, we may take $G$ to be
$\Delta$-regular. The marked edges are the original $m$ clique edges.
A \emph{pattern} on the positions $0,\ldots,n-1$ assigns one of the
colors $0,1,2$ to each unordered pair of distinct positions.
A map from these positions to $V(G)$ determines a pattern. A pair
receives color $2$ if its images are joined by a marked edge, color
$1$ if they are joined by an unmarked edge, and color $0$ if they
coincide or are nonadjacent. For instance, in
\textbf{Figure~\ref{fig:marked-maps}}(b) the two positions $0$ and
$2$ have the same image.

\begin{figure}[htbp]
\centering
\begin{tikzpicture}[every node/.style={font=\small}]
\node at (1.65,0.83) {(a) A joining edge};
\node[graph vertex,label=below:$a$] (a) at (0,0) {};
\node[graph vertex,label=below:$b$] (b) at (1.1,0) {};
\node[graph vertex,label=below:$c$] (c) at (2.2,0) {};
\node[graph vertex,label=below:$d$] (d) at (3.3,0) {};
\draw[marked edge] (a)--node[above] {$2$} (b);
\draw[unmarked edge] (b)--node[above] {$1$} (c);
\draw[marked edge] (c)--node[above] {$2$} (d);

\begin{scope}[xshift=5.2cm]
\node at (2.48,0.83) {(b) Repeated images};
\node[graph vertex,label=below:$0$] (p0) at (0,0) {};
\node[graph vertex,label=below:$1$] (p1) at (0.94,0) {};
\node[graph vertex,label=below:$2$] (p2) at (1.88,0) {};
\node[graph vertex,label=below:$u$] (u) at (4,0) {};
\node[graph vertex,label=below:$v$] (v) at (4.94,0) {};
\draw[marked edge] (p0)--node[above] {$2$} (p1);
\draw[marked edge] (p1)--node[above] {$2$} (p2);
\draw[-{Latex[length=2mm]},semithick] (2.42,0)--node[above] {$f$} (3.46,0);
\draw[marked edge] (u)--node[above] {$2$} (v);
\end{scope}
\end{tikzpicture}
\caption{The edge $bc$ joins the marked edges in (a).
In (b), $f(0)=f(2)=u$ and $f(1)=v$.}
\label{fig:marked-maps}
\end{figure}

The graph underlying a pattern $F$ has an edge for each pair of
color $1$ or $2$. For the simple-graph calculation we use patterns
in which this graph is connected, the edges of color $2$ form a strong
clique in it, and the pair $\{0,1\}$ has color $2$. We call $(0,1)$
the ordered root. Let $h_F$ be the number of maps that give exactly
the colors of $F$. To see why we use the normalized count
$q_F=\frac{h_F}{2m\Delta^{n-2}}$, fix a spanning tree of the underlying
graph that contains the root edge. Choose the image of $(0,1)$
uniformly from the $2m$ ordered marked edges. Work outwards from the
root edge along the tree. Whenever the image of a position is known,
choose the image of each of its children independently and uniformly
from the $\Delta$ neighbors of that image. Each map counted by $h_F$
occurs with probability $\frac{1}{2m\Delta^{n-2}}$. Hence $q_F$ is the
probability that the sampled map gives the prescribed colors on
\emph{all} pairs of positions.

Fix such a pattern $T$ on $k$ positions. A flag $F$ of type $T$
is one of the patterns just described whose restriction to its first
$k$ positions is $T$. Given a map $f$ that gives the colors of $T$,
let $h_F(f)$ count the maps that give the colors of $F$ and agree
with $f$ on these $k$ positions. For a flag on $n$ positions, extend
a spanning tree of $T$ to a spanning tree of $F$. Starting with the
images fixed by $f$, sample the remaining $n-k$ images along the
tree. The probability of obtaining all the colors of $F$ is now
$\frac{h_F(f)}{\Delta^{n-k}}$. This plays the role of the flag
density for a fixed placement of the labels in
\textbf{Subsection~\ref{subsec:flag-algebra}}.

Take flags $F_i$ and $F_j$ of type $T$, with disjoint sets of positions
outside $T$. An extension of $f$ to $F_i$ and an extension of $f$ to
$F_j$ together define a map on the union of their positions. Its images
specify the colors of the remaining pairs, each with one position
in $F_i$ outside $T$ and one in $F_j$ outside $T$. Let
$\mathcal U_{ij}$ be the set of patterns on the union that agree with
both flags and in which the edges of color $2$ form a strong clique.
Every pair of extensions determines exactly one pattern in
$\mathcal U_{ij}$. Conversely, a map giving the colors of any
$H\in\mathcal U_{ij}$ and extending $f$ restricts to one extension
for each flag. Therefore
\[
 h_{F_i}(f)h_{F_j}(f)=\sum_{H\in\mathcal U_{ij}}h_H(f).
\]
The identity holds even when positions outside $T$ in the two flags
have the same image, since the map records their pair as color $0$.
Thus the formal sum $\sum_{H\in\mathcal U_{ij}}H$ in
\eqref{eq:colored-flag-product} gives an exact product rule for these
counts at every $\Delta$. Freedman, Lov\'asz and Schrijver use the
same multiplication of extension counts for partially labeled
graphs~\cite{FreedmanLovaszSchrijver2007}.

In \textbf{Figure~\ref{fig:flag-union}}, only the color of $\{2,3\}$
remains to be determined. Let $H_a$ be the pattern with color $a$
on that pair. All three choices are possible in the formal product.
When $a=2$, for example, the unmarked edge on $\{0,2\}$ joins the
two disjoint marked edges on $\{0,1\}$ and $\{2,3\}$. For these flags
the sum consists of $H_0,H_1,H_2$.

\begin{figure}[htbp]
\centering
\begin{tikzpicture}[every node/.style={font=\small}]
\node[graph vertex,label=below:$0$] (a0) at (0,0) {};
\node[graph vertex,label=below:$1$] (a1) at (1.9,0) {};
\node[graph vertex,label=above:$2$] (a2) at (0,1.18) {};
\draw[marked edge] (a0)--node[above] {$2$} (a1);
\draw[unmarked edge] (a0)--node[left] {$1$} (a2);
\node at (0.95,-0.65) {First flag};

\begin{scope}[xshift=3.6cm]
\node[graph vertex,label=below:$0$] (b0) at (0,0) {};
\node[graph vertex,label=below:$1$] (b1) at (1.9,0) {};
\node[graph vertex,label=above:$3$] (b3) at (1.9,1.18) {};
\draw[marked edge] (b0)--node[above] {$2$} (b1);
\draw[unmarked edge] (b1)--node[right] {$1$} (b3);
\node at (0.95,-0.65) {Second flag};
\end{scope}

\begin{scope}[xshift=7.45cm]
\node[graph vertex,label=below:$0$] (c0) at (0,0) {};
\node[graph vertex,label=below:$1$] (c1) at (1.9,0) {};
\node[graph vertex,label=above:$2$] (c2) at (0,1.18) {};
\node[graph vertex,label=above:$3$] (c3) at (1.9,1.18) {};
\draw[marked edge] (c0)--node[above] {$2$} (c1);
\draw[unmarked edge] (c0)--node[left] {$1$} (c2);
\draw[unmarked edge] (c1)--node[right] {$1$} (c3);
\draw[densely dashed] (c2)--node[above] {$?$} (c3);
\node at (0.95,-0.65) {Union};
\end{scope}
\end{tikzpicture}
\caption{Two flags with common type on $\{0,1\}$. In their union,
the dashed pair has color $0$, $1$ or $2$. Undrawn pairs have color $0$.}
\label{fig:flag-union}
\end{figure}

The product rule gives inequalities as follows. Fix a list of flags of
type $T$. For the $i$th flag $F_i$, on $n_i$ positions, put
$v_i(f)=\frac{h_{F_i}(f)}{\Delta^{n_i-k}}$. Let $v(f)$ be the column
vector with these entries. The matrix
\[
 \frac{1}{2m\Delta^{k-2}}\sum_f v(f)v(f)^{\mathsf T}
\]
is positive semidefinite, since each summand is an outer product.
The sum ranges over the maps giving the colors of $T$. In its
$(i,j)$ entry the product rule sums the counts of patterns
$H\in\mathcal U_{ij}$ on $n_i+n_j-k$ positions. Each map counted by
$h_H$ has a unique restriction to $T$, so $\sum_f h_H(f)=h_H$.
The three factors in the denominator give
$2m\Delta^{k-2}\Delta^{n_i-k}\Delta^{n_j-k}
=2m\Delta^{n_i+n_j-k-2}$. Thus the entry is
$\sum_{H\in\mathcal U_{ij}}q_H$, as in \eqref{eq:moment-expansion}.
In the example of \textbf{Figure~\ref{fig:flag-union}} it is
$q_{H_0}+q_{H_1}+q_{H_2}$. Consequently each quadratic form of this matrix, evaluated at a fixed
real vector, gives a linear expression in the counts $q_H$ whose
value is nonnegative.

For a multigraph, a map into the simple support records which pairs
of vertices are adjacent. We also have to record how many edges of
the strong clique join each pair. Call an edge of the support
\emph{red} if at least one marked copy joins its endpoints. For each
red edge $uv$, distinguish one marked copy, called the \emph{base
copy}, and let $b_{uv}$ be the number of other marked copies, called
\emph{extra copies}. A pattern can now require an extra copy on one
or more specified edges. If the map sends such an edge to $uv$, there
are $b_{uv}$ choices for that copy. For a fixed map, multiplying
these numbers counts the choices on all specified edges.

For two flags with a common type, first choose the copies required
on the type. The choices on each flag outside the type are then
independent, so the same product rule counts pairs of extensions.
When the type itself requires extra copies, its map $f$ contributes
an outer product multiplied by the number of choices on the type.
This number is nonnegative, and the resulting matrix is again positive
semidefinite. \textbf{Subsection~\ref{sec:weighted-counts}} specifies
the weighted counts used in the multigraph proof.

\subsection{The Uniform Bound for Simple Graphs}\label{subsec:marked-edge-pairs}

We now express $\frac{m}{\Delta^2}$ in terms of the counts $q_F$.
Choose the images of $(0,1)$ and $(2,3)$ independently from the $2m$
ordered marked edges. The two choices may give the same edge. Each
choice determines the images of four positions and hence the colors of
all six pairs. If the two image edges are disjoint, a joining edge
exists because the marked edges form a strong clique. If they share
an endpoint, a position for that endpoint is adjacent in the pattern
to a position on the other edge. Thus the pattern is connected in
either case. Let $\mathcal B$ consist of the connected patterns on
these four positions whose edges of color $2$ form a strong clique
and include $\{0,1\}$ and $\{2,3\}$. Each choice of the two ordered
edges contributes to exactly one $h_F$ with $F\in\mathcal B$, and
every map counted by such an $h_F$ determines the two ordered edges.
Therefore
\[
 \sum_{F\in\mathcal B}h_F=(2m)^2,\qquad
 \frac{1}{2}\sum_{F\in\mathcal B}q_F=\frac{m}{\Delta^2}.
\]

To bound this sum we use the flag inequalities from
\textbf{Subsection~\ref{subsec:adapted-flags}} and the degree
equations. For the latter, fix a position in a pattern and require
one new position to map to a neighbor of its image. Every map of the
old positions has exactly $\Delta$ extensions. Sorting these
extensions by their pair colors gives \eqref{eq:degree-extension}
after normalization. The calculation in
\textbf{Subsection~\ref{sec:simple-proof}} uses the degree equations to
express $\frac{5}{4}-\frac{1}{2}\sum_{F\in\mathcal B}q_F$ as a sum
of nonnegative multiples of individual counts $q_F$ and of quadratic
forms of the positive semidefinite matrices with nonnegative weights. The identity is \eqref{eq:certificate-identity},
with rational coefficients checked in the proof of
\textbf{Lemma~\ref{prop:checked-certificate}}. Its right-hand side
is nonnegative. The identity for the two ordered edges now gives
$m\leq\frac{5}{4}\Delta^2$.

\subsection{From Simple Graphs to Multigraphs}\label{subsec:weighted-sketch}

Let $D$ be a positive integer and let $G$ be a loopless multigraph of
maximum degree at most $D$. Mark the $m>0$ edges of a strong clique.
Keep the subgraph induced by their endpoints. Between vertices joined
by marked edges, keep all the marked copies. Between any other two
adjacent vertices, keep one unmarked edge. Thus each pair of disjoint
marked edges still has a joining edge. Let $s$ be the number of red
edges in the simple support and $b=m-s$ the total number of extra
copies. The $s$ red edges form a strong clique in the simple support,
so the simple-graph bound gives $s\leq\frac{5}{4}D^2$.

The degree of a vertex in the reduced multigraph is its number of
distinct neighbors plus the number of extra copies incident with it.
As shown in \textbf{Subsection~\ref{sec:weighted-counts}}, we may add
unmarked edges and vertices to make the multigraph $D$-regular while
preserving its induced subgraph on the old vertices. Then, for each
vertex $v$, the extension by one incident edge has exactly $D$
choices. A neighbor $u$ joined to $v$ by an unmarked edge contributes
one choice. A neighbor on a red edge $uv$ contributes one choice of
the base copy and $b_{uv}$ choices of an extra copy. Counting these
choices in the weighted patterns gives the degree equations for
multigraphs.

Choose two ordered marked copies as in the simple-graph calculation.
There are $2s$ choices for an ordered base copy and $2b$ for an ordered
extra copy. Thus the numbers of choices using zero, one and two extra
copies are $(2s)^2$, $2(2s)(2b)$ and $(2b)^2$, respectively. Dividing
each count by $4mD^2$ gives the identity
\[
 \frac{s^2}{mD^2}+\frac{2sb}{mD^2}+\frac{b^2}{mD^2}
 =\frac{m}{D^2}.
\]
Each term on the left is a linear combination of the weighted
pattern counts. For example, in the last term each chosen edge of
the support contributes its number of extra copies.

Write $M=\frac{m}{D^2}$, $S=\frac{s}{D^2}$ and $e=\frac{b}{m}$.
Then $S=M(1-e)$. Suppose that $M>\frac{5}{4}$. The weighted flag
inequalities first give upper bounds on both $M$ and $e$. Degree
estimates on induced five-cycles of red edges, together with flag
inequalities in the simple support, give an upper bound on $S$.
The identity $e=1-\frac{S}{M}$ then gives a lower bound on $e$.
With $e$ restricted to this interval, the weighted flag inequalities
give a smaller upper bound on $M$. The successive bounds in the proof
of \textbf{Lemma~\ref{prop:multi-uniform}} yield $M<\frac{5}{4}$,
contradicting our assumption. This proves $m\leq\frac{5}{4}D^2$.
For even $D$, the right-hand side is $f(D)$.

\subsection{The Structure of a Counterexample}\label{subsec:odd-sketch}

Let $D=\Delta(G)$ be odd. If the bound $f(D)$ fails, mark exactly
$m=f(D)+1$ edges of a strong clique and keep the multigraph on their
endpoints as in \textbf{Subsection~\ref{subsec:weighted-sketch}}.
The definition of $f(D)$ gives
\[
 M=\frac{m}{D^2}=\frac{5}{4}-\frac{1}{2D}+\frac{5}{4D^2}.
\]
For $D\geq9$ this value is at least $\frac{98}{81}$, with equality
at $D=9$. If the simple support contained no induced five-cycle with
all five edges red, the flag inequality in~\cite[(G.2)]{FGSTSupplement} would give
$M<\frac{98}{81}$. Thus the support contains such a cycle, as stated
in \textbf{Lemma~\ref{lem:counterexample-cycle}}. We examine the degrees around
each such cycle to bound $s$, the number of distinct pairs of
endpoints of the $m$ marked edges.

Fix one of these cycles, say $c_0c_1c_2c_3c_4c_0$, and take
indices modulo five. For each $i$, let $U_i$ consist of the vertices
whose only neighbors on the cycle are $c_{i-1}$ and $c_{i+1}$ and
whose edges to both of these vertices are red. In particular,
$c_i\in U_i$. Take $x\in U_0$ and $y\in U_1$. The disjoint red edges
$xc_4$ and $yc_2$ must have a joining edge. The definitions of
$U_0,U_1$ exclude $xc_2$ and $yc_4$, and the induced cycle has no
edge $c_4c_2$. Thus $xy$ is an edge of the support. The same argument
with indices shifted shows that each vertex of $U_i$ is adjacent to
every vertex of $U_{i-1}$ and $U_{i+1}$.

In particular, $c_i$ is joined to at least $|U_{i-1}|+|U_{i+1}|$
distinct vertices in the multigraph. After counting one copy of each
of these edges, at most $D-|U_{i-1}|-|U_{i+1}|$ more edges can be
incident with $c_i$. The sum of these five upper bounds is
$R=5D-2\sum_{i=0}^4|U_i|$. For odd $D\geq11$, the degree count in
\textbf{Subsection~\ref{subsec:exact-large-degree}} shows that $R>\frac{D+1}{2}$ on
every such cycle in a counterexample.

To see how this structural restriction yields an inequality, let
$q_C$ be the number of ordered red induced five-cycles divided by
$2sD^3$. Here an ordered cycle has a specified starting vertex and
direction. Let $L_R$ be the sum of $\frac{R}{D}$ over the same
ordered cycles, divided by the same denominator. The bound on $R$
gives $L_R\geq\frac{q_C}{2}$. One of the flag inequalities for the
simple support in~\cite[(G.1)]{FGSTSupplement} gives
\[
 S+0.45119\left(L_R-\frac{q_C}{2}\right)\leq 1.221800627,
 \qquad\text{where }S=\frac{s}{D^2}.
\]
Its second term on the left is nonnegative. Thus
$S\leq1.221800627$. For odd $D\geq15$, the value of $M$ above
satisfies $M\geq\frac{11}{9}>1.221800627$. It follows that $s<m$,
so some pair of vertices is joined by at least two marked copies.
Moreover, $e=1-\frac{S}{M}$ gives an explicit lower bound on the
fraction of extra copies.

For $D=11,13$, the integer $R$ is odd, and the degree estimates
in \textbf{Subsection~\ref{subsec:exact-large-degree}} give $R\geq9$. Substituting
$L_R\geq\frac{9}{D}q_C$ into the flag inequalities gives $S<M$
in these two cases as well. At $D=9$, the finite simple-graph
bound excludes $s=m$. The remaining degree estimates restrict $e$
further. For each odd $D\geq9$ the weighted flag inequalities then
exclude the prescribed value $M=\frac{f(D)+1}{D^2}$ throughout the
allowed interval for $e$, as detailed in
\textbf{Table~\ref{tab:multi-odd}}. The cases $D=1,3,5,7$ are
settled separately in \textbf{Subsection~\ref{subsec:exact-small-degree}}. This proves
\textbf{Theorem~\ref{thm:multigraph}}.

